\documentclass[10pt,a4paper]{amsart}
\usepackage[margin=19mm]{geometry}
\usepackage{amsmath,amssymb,mathtools}
\usepackage[scr=rsfs,cal=euler]{mathalfa}
\usepackage{xfrac}
\usepackage[unicode,hidelinks,bookmarks=false]{hyperref}
\newcommand{\ie}{i.e.\ }
\newcommand{\cf}{cf.\ }
\newcommand{\resp}{respectively\ }
\newcommand{\Subsection}{\subsection}
\providecommand{\nobreakdash}{\nobreak\hbox{-}}
\setlength{\emergencystretch}{2em}
\multlinegap0pt
\usepackage{bm}
\usepackage{bbm}
\usepackage{dsfont}
\usepackage{xspace}
\usepackage{cancel}
\usepackage{mathtools}
\usepackage{amsthm}
\makeatletter
\@ifundefined{theorem}{\newtheorem{theorem}{Theorem}}{}
\@ifundefined{lemma}{\newtheorem{lemma}[theorem]{Lemma}}{}
\makeatother
\newcommand{\dx}{\textup{d}x}
\title[Global weak solutions to the Cahn-Hilliard equation with deg]{Global weak solutions to the Cahn-Hilliard equation with degenerate mobility and singular diffusion}
\author{Monica Conti, Andrea Giorgini, Greta Ricchi}
\date{Source: arXiv:2608.06371v2 (CC BY 4.0)}
\begin{document}
\maketitle

\noindent\emph{Source and licence.} Global weak solutions to the Cahn-Hilliard equation with degenerate mobility and singular diffusion. Version arXiv:2608.06371v2 (CC BY 4.0), distributed under the Creative Commons Attribution 4.0 International licence, as recorded in the arXiv licence metadata for this exact version and shown on the abstract page https://arxiv.org/abs/2608.06371v2. Full rights notice: licenses/arxiv-2608.06371-CC-BY-4.0.txt.\par\smallskip

\noindent\textit{Editorial scope.} Counted unit: the Lemma whose source label is \texttt{refined-lemma} in the article (source0.tex lines 738--765, source label \texttt{refined-lemma}), delivered together with the article's own complete original proof of it (source0.tex lines 767--966, one \texttt{proof} environment of 200 lines). No mathematics is written, repaired or re-derived: statement and proof are the article's own text line for line. Required context retained uncounted: lemma H2 (lines 513--521). Every definition, display and label of those lines is retained unchanged. Omitted from this package and not redistributed: 1--512, 522--737, 766--766, 967--2136. Nothing inside a retained excerpt is omitted or altered: every retained body is delivered exactly as the article's lines are written, so no omission receipt of any kind accompanies this package. Citations: the retained excerpts carry no citation command at all, so no bibliography entry of the article is transcribed and no reference is redistributed. Reference adaptation: none was needed and none was made; every reference inside the delivered unit resolves inside it, so no unresolved reference or citation is produced. Printed slips kept literal: no printed word, symbol, hypothesis or number of the counted statement or of its proof was altered. Layout and class: the article's own class is \texttt{amsart} with packages including fullpage, amsmath, amsthm, mathrsfs, mathtools, amsfonts, amssymb, amsaddr, tikz-cd, inputenc, enumerate, babel, esint, hyperref; the wrapper is the lane's pinned \texttt{amsart} editorial wrapper at 10pt on A4, which restates the theorem-like environments the retained text needs and adds the article's own macro definitions that the retained text needs (\texttt{\textbackslash dx}), with extra wrapper packages (bm, bbm, dsfont, xspace, cancel, mathtools). The delivered numbering is the unit's own. Assets: no retained span contains a figure environment, a graphics-inclusion command, a PDF-page inclusion, a table or a TikZ picture; no image file is copied into this package, the package declares no asset dependency and no image file is redistributed with it. Licence: version arXiv:2608.06371v2 of this article is distributed under the Creative Commons Attribution 4.0 International licence, as recorded in the arXiv licence metadata for this exact version and shown on the abstract page https://arxiv.org/abs/2608.06371v2. A full rights notice is delivered at licenses/arxiv-2608.06371-CC-BY-4.0.txt. Characters: every retained excerpt is pure ASCII, so the delivered engine, XeLaTeX with its default Latin Modern text font, reports no missing character.
\begin{lemma}
\label{lemma H2}
    Let $\Omega$ be a smooth convex set. Assume that
    $f \in W^{2,2}(\Omega)$ with $\partial_{\mathbf{n}} f=0$ on $\partial \Omega$. Then
    $$
    \int_\Omega |\Delta f|^2 \ \dx \ge \int_\Omega |D^2 f |_F^2 \ \dx,
    $$
where $D^2 f$ is the Hessian of $f$ and $|A|_F= \sqrt{\textrm{tr}(A^TA)}$ is the Frobenius norm of a square matrix $A$. 
\end{lemma}

\begin{lemma}
\label{refined-lemma}
Let $\Omega\subset\mathbbm{R}^d$ be a bounded, smooth and convex domain with $d=2,3$, and let $v\in H^2(\Omega)$ satisfy $\partial_\mathbf{n} v=0$ on $\partial \Omega$.
Let $J\subset\mathbbm{R}$ be an interval such that $v(x)\in J$ for all
$x\in\overline{\Omega}$, and assume that $H\in W^{2,\infty}(J)$. Set
\[
    I_H
    :=
    \int_\Omega H'(v)|\nabla v|^2\Delta v \ \dx.
\]
Then,
\begin{align}
    I_H+\int_\Omega|\Delta v|^2\ \dx
    \geq{}&
    \frac1d\int_\Omega|\Delta v|^2\ \dx
    \nonumber\\
    &+
    \frac{d}{d+2}
    \int_\Omega
    \left[
        -H''(v)
        -
        \frac{d}{d+2}|H'(v)|^2
    \right]
    |\nabla v|^4\ \dx.
\label{eq:refined-lemma}
\end{align}
\end{lemma}

\begin{proof}
By density it is sufficient to assume that $v \in H^3(\Omega)$ with 
$v(x)\in J$ for all $x\in\overline{\Omega}$ and
$\partial_{\mathbf{n}} v=0$ on $\partial \Omega$.
We first derive an identity for $I_H$. Consider the vector field
\[
    X:=H'(v)|\nabla v|^2\nabla v.
\]
By the product rule,
\begin{align*}
    \operatorname{div}X
    &=
    H''(v)|\nabla v|^4
    +
    H'(v)\nabla|\nabla v|^2\cdot\nabla v
    +
    H'(v)|\nabla v|^2\Delta v.
\end{align*}
Since
\[
    \frac12\nabla|\nabla v|^2=D^2v\,\nabla v,
\]
we obtain
\[
    \operatorname{div}X
    =
    H''(v)|\nabla v|^4
    +
    2H'(v)  D^2v\,\nabla v \cdot \nabla v
    +
    H'(v)|\nabla v|^2\Delta v.
\]
Integrating over $\Omega$ and using the divergence theorem gives
\[
    \int_\Omega\operatorname{div}X\ \dx
    =
    \int_{\partial\Omega}
    H'(v)|\nabla v|^2\partial_{\mathbf{n}} v\ 
    \mathrm{d}\mathcal H^{d-1}
    =
    0,
\]
where the last equality follows from the homogeneous Neumann boundary
condition. Therefore,
\begin{equation}
\label{eq:A-identity}
    I_H
    =
    -\int_\Omega H''(v)|\nabla v|^4\ \dx
    -
    2\int_\Omega
    H'(v)
     D^2v\,\nabla v \cdot \nabla v \ \dx.
\end{equation}
Adding $\int_\Omega|\Delta v|^2\ \dx$ to both sides and using Lemma \ref{lemma H2}, we have
\begin{align}
    I_H+\int_\Omega|\Delta v|^2\ \dx
    \geq
    \int_\Omega
    \Big[
        |D^2v|_F^2
        -
        2H'(v)
         D^2v\,\nabla v \cdot \nabla v
        -H''(v)|\nabla v|^4
    \Big] \ \dx.
\label{eq:starting-cances}
\end{align}

We now use once again \eqref{eq:A-identity}, equivalently written as
\begin{equation}
\label{eq:zero-identity}
    0
    =
    \int_\Omega
    \Big[
        H''(v)|\nabla v|^4
        +
        2H'(v)
        D^2v\,\nabla v \cdot \nabla v
        +
        H'(v)|\nabla v|^2\Delta v
    \Big]\ \dx.
\end{equation}
Multiplying \eqref{eq:zero-identity} by $2/(d+2)$ and adding it to
the right-hand side of \eqref{eq:starting-cances}, we obtain
\begin{align}
    I_H+\int_\Omega|\Delta v|^2\ \dx
    \geq
    \int_\Omega
    \Big[
        |D^2v|_F^2
        &-
        \frac{2d}{d+2}
        H'(v)
         D^2v\,\nabla v \cdot \nabla v 
        \nonumber\\
        &+
        \frac{2}{d+2}
        H'(v)|\nabla v|^2\Delta v
        -
        \frac{d}{d+2}
        H''(v)|\nabla v|^4
    \Big] \ \dx.
\label{eq:cances-combination}
\end{align}

Define the matrix
\[
    \mathbb{D}(v)
    :=
    D^2v-\frac{\Delta v}{d}\mathbbm{I}_d,
\]
where $\mathbb{I}_d$ is the identity matrix.
Since the matrix $D^2 v$ is symmetric,
%$$\nabla^2 \phi_\ve(u_\ve)= \mathbf{R}+\dfrac{\Delta \phi_\ve(u_\ve)}{3} \mathbb{I}_3$$
and 
$\Delta v= \textrm{tr}(D^2 v)$ leads to $\textrm{tr}\left(  \mathbb{D}(v) \right)=0$, we infer that
\begin{equation}
\label{eq:R-norm}
\left|D^2 v \right|_F^2
= \textrm{tr}\left[ \left(  \mathbb{D}(v)  +\dfrac{\Delta v}{d} \mathbbm{I}_d\right)^2\right]
= \left|  \mathbb{D}(v)  \right|_F^2+\dfrac{|\Delta v|^2}{d}
\end{equation}
and
\begin{equation}
\label{eq:R-mixed}
    D^2v\,\nabla v \cdot \nabla v
    =
  \mathbb{D}(v) \, \nabla v \cdot \nabla v
    +
    \frac{\Delta v}{d}|\nabla v|^2.
\end{equation}
Substituting \eqref{eq:R-norm} and \eqref{eq:R-mixed} into
\eqref{eq:cances-combination}, the mixed terms involving
$H'(v)|\nabla v|^2\Delta v$ cancel exactly. Hence, we arrive at
\begin{align}
    I_H+\int_\Omega|\Delta v|^2\ \dx
    & \geq \
    \frac1d \int_\Omega|\Delta v|^2\ \dx
    \nonumber\\
    &\quad +
    \int_\Omega
    \Big[
        | \mathbb{D}(v) |_F^2
        -
        \frac{2d}{d+2}
        H'(v)   \mathbb{D}(v) \nabla v \cdot \nabla v
        -
        \frac{d}{d+2}
        H''(v)|\nabla v|^4
    \Big]\ \dx.
\label{eq:trace-free}
\end{align}

Set
\[
    a:=\frac{d}{d+2}.
\]
The last integrand can be completed to a square as follows
\begin{align}
    &| \mathbb{D}(v) |_F^2
    -
    2a H'(v)  \mathbb{D}(v) \nabla v \cdot \nabla v
    -
    aH''(v)|\nabla v|^4
    \nonumber\\
    &\qquad=
    \left|
         \mathbb{D}(v) -aH'(v)\nabla v\otimes\nabla v
    \right|_F^2
    +
    a
    \Big[
        -H''(v)-a|H'(v)|^2
    \Big]
    |\nabla v|^4.
\label{eq:square-completion}
\end{align}
Since the first term on the right-hand side of
\eqref{eq:square-completion} is non-negative, we infer from
\eqref{eq:trace-free} that
\begin{align}
    I_H+\int_\Omega|\Delta v|^2\ \dx
    &\geq
    \frac1d \int_\Omega|\Delta v|^2\ \dx
    \nonumber\\
    &\quad 
    +
    \frac{d}{d+2}
    \int_\Omega
    \left[
        -H''(v)
        -
        \frac{d}{d+2}|H'(v)|^2
    \right]
    |\nabla v|^4\ \dx,
\end{align}
which proves \eqref{eq:refined-lemma}.
\end{proof}

\end{document}
